\documentclass[11pt,a4paper]{article}
\usepackage[T1]{fontenc}
\usepackage[utf8]{inputenc}
\usepackage[polish]{babel}
\usepackage{polski}
\usepackage{geometry}
\usepackage{enumitem}
\usepackage{hyperref}
\usepackage{xcolor}
\usepackage{fancyhdr}
\usepackage{parskip}

\geometry{margin=2.2cm}
\hypersetup{colorlinks=true,linkcolor=blue!50!black,urlcolor=blue!60!black}
\pagestyle{fancy}
\fancyhf{}
\fancyhead[L]{\small Ideo -- Flutter}
\fancyhead[R]{\small Adrian -- proste Q\&A}
\fancyfoot[C]{\thepage}
\renewcommand{\headrulewidth}{0.4pt}

\newcommand{\pyt}[1]{\par\medskip\noindent\textbf{Pytanie:} #1\par}
\newcommand{\odp}[1]{\noindent\textbf{Odpowiedź:} #1\par}

\title{Rozmowa w Ideo -- Flutter\\[0.4em]\large Proste pytania i odpowiedzi}
\author{Adrian Rojek\\11.09.2026, godz.\ 10:00\\ul.\ Nad Przyrwą 13}
\date{}

\begin{document}
\maketitle
\tableofcontents
\newpage

\section{Czy będą kazać kodować na żywo?}

\textbf{Raczej nie.} Spotkanie trwa max 45 minut. Będą 2 osoby techniczne. Więcej czasu zajmie gadanie o Tobie, projektach i firmie niż siedzenie przy laptopie.

Przygotuj się na:
\begin{itemize}
  \item pytania o Fluttera, Firebase, REST, stan aplikacji,
  \item opowiadanie o swoich apkach,
  \item ewentualnie krótkie ,,jak byś to zrobił'' na tablicy / ustnie.
\end{itemize}
Jak będzie druga runda -- dopiero wtedy częściej bywa zadanie albo live coding.

\section{Powiedz o sobie w 30 sekund}

\pyt{Opowiedz coś o sobie.}
\odp{Jestem studentem informatyki w Rzeszowie. Robię aplikacje mobilne -- głównie Flutter i SwiftUI. W DEVEJI robiłem Fluttera z Firebase i wrzucałem apki do App Store i Google Play. Mam też własną apkę Flutter (owoce/warzywa, skan ML, SQLite) oraz apki iOS typu Water Reminder i Fitness Tracker.}

\pyt{Dlaczego Ideo?}
\odp{Szukacie Fluttera, Firebase, REST, ładnego UI i kogoś kto ogarnia też natywne iOS/Android. To pasuje do tego, co już robiłem. Chcę robić realne projekty dla klientów w software house i uczyć się od zespołu. Rzeszów i hybryda mi pasują.}

\section{Pytania miękkie / o pracę}

\pyt{Jak wyglądała praca w DEVEJI?}
\odp{Robiłem ekrany we Flutterze, łączyłem z Firebase, poprawiałem bugi, przygotowywałem buildy do sklepów. Jak coś było niejasne albo ryzykowne (np.\ review w App Store) -- mówię wcześniej, nie na ostatnią chwilę.}

\pyt{Co robisz, gdy wymagania są niejasne?}
\odp{Dopytuję. Ustalam najprostsze MVP. Wolę potwierdzić założenie niż zgadywać przez dwa dni.}

\pyt{Co gdy utkniesz na błędzie?}
\odp{Odtwarzam problem na minimum. Patrzą logi. Odcinam warstwy: UI / stan / sieć. Jak blokuję innych -- mówię szybko i proponuję plan B.}

\pyt{Remote, hybryda czy biuro?}
\odp{Hybryda. Na start chętnie bywam w biurze, potem według projektu.}

\pyt{Masz apki w sklepach?}
\odp{Tak -- przy DEVEJI brałem udział w publikacji iOS i Androida. Mam też własne projekty iOS. O statusie Fruit Trackera mów tylko zgodnie z prawdą.}

\pyt{Dasz radę też przy innych technologiach?}
\odp{Tak. Flutter będzie główny, ale znam SwiftUI, trochę Kotlina, robiłem też Electron/Node i proste strony. Oferta Ideo właśnie to przewiduje.}

\pyt{Studia + praca?}
\odp{Robię magistra. Potrafię łączyć -- sesję i wolne komunikuję z wyprzedzeniem.}

\section{Flutter i Dart -- prosto}

\pyt{Co to Flutter?}
\odp{Narzędzie do robienia aplikacji mobilnych. Piszesz w Dart. Jedna baza kodu na Androida i iOS. Ekrany składasz z widgetów.}

\pyt{Stateless vs Stateful?}
\odp{Stateless -- ekran bez własnego zmieniającego się stanu. Stateful -- ma lokalny stan i \texttt{setState}. Do większej logiki lepiej Provider/Cubit/BloC.}

\pyt{Co to BuildContext?}
\odp{To ,,adres'' widgetu w drzewie. Dzięki niemu bierzesz motyw, nawigację, Provider. Po \texttt{await} sprawdzaj \texttt{context.mounted}, zanim użyjesz contextu.}

\pyt{Po co Keys?}
\odp{Żeby Flutter dobrze rozpoznawał elementy listy przy zmianach kolejności i nie mieszał stanów.}

\pyt{Po co const?}
\odp{Mniej zbędnego przebudowywania UI. Gdzie się da -- dawaj \texttt{const}.}

\pyt{async / await?}
\odp{Na rzeczy które trwają: sieć, baza. UI działa dalej. Ciężkie liczenie lepiej w osobnym isolate (\texttt{compute}).}

\pyt{Future vs Stream?}
\odp{Future = jeden wynik. Stream = wiele wyników w czasie (np.\ dane z bazy na żywo).}

\pyt{Null safety?}
\odp{Zmienna domyślnie nie może być null. Jak może -- dajesz \texttt{?}. \texttt{!} używaj tylko gdy naprawdę wiesz, że nie ma nulla.}

\pyt{Hot reload vs hot restart?}
\odp{Reload -- szybka zmiana UI ze stanem. Restart -- od zera. Pełny restart aplikacji przy zmianach natywnych.}

\pyt{ListView vs ListView.builder?}
\odp{Zwykły ListView buduje wszystko naraz -- źle przy długiej liście. Builder buduje tylko to, co widać.}

\pyt{Material vs Cupertino?}
\odp{Material = styl Androida/Google. Cupertino = styl iOS. Często robimy według designu klienta.}

\section{Stan aplikacji: Provider, Cubit, BloC}

\pyt{Jakie znasz sposoby na stan?}
\odp{setState lokalnie. Provider. Cubit. BloC. W ofercie Ideo są Cubit, Provider i BloC -- o tych mów śmiało. W Fruit Trackerze używałem Providera.}

\pyt{Provider -- co to?}
\odp{Trzymasz stan w klasie (np.\ ChangeNotifier) i UI go słucha. Proste i dobre na średnie apki. Nie rób jednego wielkiego obiektu na wszystko.}

\pyt{Cubit vs BloC?}
\odp{Cubit: wołasz metodę i emitujesz nowy stan -- mniej kodu. BloC: wysyłasz eventy, dostajesz stany -- jaśniejszy przepływ przy trudniejszej logice. Mały feature = często Cubit. Duży, skomplikowany = BloC.}

\pyt{Kiedy BloC zamiast Providera?}
\odp{Gdy logika jest grubsza, chcesz stany loading/error/data, łatwe testy i czytelny przepływ. Provider wystarczy na prostszy stan UI.}

\pyt{Jak uniknąć przebudowy całego ekranu?}
\odp{Słuchaj tylko potrzebnego kawałka stanu. Dziel providery/cubity. Stałe części UI dawaj jako \texttt{const}.}

\pyt{Gdzie robić nawigację i snackbary?}
\odp{Nie w metodzie \texttt{build}. Lepiej w listenerze po zmianie stanu albo po skończonym Future.}

\section{REST -- rozmowa z serwerem}

\pyt{Jak wołasz API we Flutterze?}
\odp{Pakiet \texttt{http} albo \texttt{dio}. Warstwy: API client $\to$ repository $\to$ Cubit/BloC $\to$ UI. Dane z JSON mapujesz na modele.}

\pyt{Błędy sieci?}
\odp{Timeout. Pokazujesz błąd w UI. Przy GET możesz spróbować ponowić. Logujesz crash/non-fatal. Offline: cache albo komunikat.}

\pyt{Logowanie / token?}
\odp{Token w nagłówku. Trzymasz w secure storage, nie w zwykłych preferencjach na produkcji. Przy 401 -- odśwież token albo wyloguj.}

\pyt{Paginacja?}
\odp{Dociągasz kolejne strony przy scrollu. Pokazujesz loader na dole. Nie dubluj tych samych elementów.}

\section{Firebase}

\pyt{Czego używałeś z Firebase?}
\odp{W DEVEJI Firebase był backendem pod apki mobilne (Auth, baza, ewentualnie push/crash -- mów według faktów z projektu). Umiem spiąć Fluttera z Firebase na iOS i Androidzie.}

\pyt{Firestore vs Realtime Database?}
\odp{Firestore = dokumenty i kolekcje (częściej dziś). Realtime = drzewo JSON, szybkie sync. Do nowych rzeczy zwykle Firestore.}

\pyt{Jak chronisz dane?}
\odp{Security Rules. Użytkownik widzi tylko swoje rzeczy po \texttt{auth.uid}. Nie ufaj samym regułom po stronie apki.}

\pyt{Auth -- jak działa w apce?}
\odp{Logowanie (mail/Google/Apple). Słuchasz czy user jest zalogowany. Na tej podstawie pokazujesz login albo home.}

\pyt{Push (FCM)?}
\odp{Token telefonu, zgody na iOS, obsługa powiadomienia gdy apka jest otwarta/w tle, przejście do właściwego ekranu.}

\pyt{Crashlytics?}
\odp{Widzisz crashe z produkcji. Bez tego trudno naprawiać problemy użytkowników.}

\section{UI i design}

\pyt{Jak dbasz o zgodność z projektem graficznym?}
\odp{Sprawdzam odstępy, fonty, kolory z Figmy. Używam theme. Porównuję na telefonie, nie tylko w emulatorze.}

\pyt{Różne rozmiary ekranów?}
\odp{Flexible/Expanded, LayoutBuilder, unikanie sztywnych wysokości na ślepo.}

\pyt{Animacje?}
\odp{Proste AnimatedContainer, Hero, czasem Lottie. Animacja ma pomagać, nie rozpraszać.}

\section{Natywne Android / iOS}

\pyt{Kiedy piszesz kod natywny przy Flutterze?}
\odp{Gdy brakuje gotowego pluginu albo trzeba głęboko wejść w system (specjalne SDK, HealthKit-like, tło). Most: MethodChannel.}

\pyt{Masz doświadczenie natywne?}
\odp{Tak, iOS/SwiftUI: Water Reminder, Fitness Tracker, Lotto. Android natywnie znam podstawowo -- w praktyce głównie Flutter + konfiguracja projektu.}

\pyt{Różnice iOS vs Android?}
\odp{Uprawnienia, push, nawigacja wstecz, zasady sklepów, różne telefony z Androidem.}

\section{Wrzucanie do sklepów}

\pyt{Jak wypuszczasz apkę?}
\odp{Podbijasz wersję, podpisujesz, budujesz AAB/IPA, uzupełniasz opisy i privacy, testujesz na prawdziwym telefonie, wrzucasz do review / staged rollout.}

\pyt{Dlaczego App Store odrzuca?}
\odp{Crash na starcie, brakujące privacy, niedziałający login, brak konta demo dla reviewera, łamanie zasad płatności.}

\pyt{AAB vs APK?}
\odp{Play chce AAB. Z niego Google robi APK pod konkretne telefony.}

\section{Architektura i git}

\pyt{Jak dzielisz kod we Flutterze?}
\odp{UI osobno, logika (Cubit) osobno, dane (API/baza) osobno. Widget nie powinien robić wszystkiego.}

\pyt{Co poprawisz w swoich starych apkach Swift?}
\odp{Za dużo logiki w jednym pliku ContentView. Wydzieliłbym ViewModel i serwisy. Dodałbym testy do ważnych reguł (np.\ streak).}

\pyt{Git?}
\odp{Branch na feature, małe PR-y, czytelne commity. Nie wrzucam haseł i \texttt{.env} do repo.}

\pyt{Testy?}
\odp{W prywatnych projektach bywało słabo -- wiem, że to minus. W pracy chcę testować logikę Cubitów i główne ścieżki. Umiesz powiedzieć jak przetestować prostą rzecz.}

\section{Twoje projekty -- co mówić}

\subsection{DEVEJI}
\pyt{Co robiłeś?}
\odp{Flutter iOS/Android, Firebase, publikacja w sklepach (VIII--X 2025). To Twój najmocniejszy punkt pod Ideo. Bez NDA: mów kategoriami zadań, nie tajnymi szczegółami klienta.}

\subsection{Fruit Tracker (Flutter, prywatne)}
\pyt{Co to?}
\odp{Apka do liczenia owoców i warzyw. Skan ML, historia, ulubione, lista zakupów offline. Stan: Provider. Dane: SQLite na telefonie.}

\pyt{Czemu SQLite a nie tylko chmura?}
\odp{Ma działać offline i trzymać dane lokalnie. Chmurę można dodać później.}

\pyt{ML -- jak to działa w skrócie?}
\odp{Zdjęcie $\to$ model zgaduje produkt $\to$ user potwierdza $\to$ zapis do historii. Przy słabym świetle model może się mylić, dlatego potwierdzenie.}

\subsection{Water Reminder}
\pyt{Co to?}
\odp{Apka iOS do picia wody. Cel 4000 ml, szybkie dodawanie, streak, powiadomienia, historia. Persistence: \textbf{SwiftData} (nie Core Data).}

\pyt{Co byś poprawił?}
\odp{Rozbić jeden duży ContentView na mniejsze części i osobną logikę.}

\subsection{Fitness Tracker}
\pyt{Co to?}
\odp{Kroki z HealthKit, przybliżone kalorie (MET), waga, wykresy, profil. Mów tylko o tym, co naprawdę jest w kodzie -- bez zmyślania MapKit/zdjęć jeśli ich nie ma.}

\subsection{MacSpaceCleaner}
\pyt{Po co o tym na rozmowie Flutter?}
\odp{Pokazuje, że myślisz o architekturze i bezpieczeństwie: osobna warstwa Core/UI, reguły co wolno usuwać, potwierdzenia. To przenosi się na mobile.}

\subsection{Courier (Electron)}
\odp{System dla firmy kurierskiej: role, zadania, tracking, PDF. Pokazuje pracę z procesem biznesowym, nie tylko ładnym UI.}

\subsection{Bus (Laravel)}
\odp{Trasy i rozkłady, role admin/kierowca. To CRUD -- \textbf{bez} prawdziwego GPS na żywo. Lepiej powiedzieć prawdę niż marketing.}

\section{Trudne pytania -- uczciwe odpowiedzi}

\pyt{Na GitHubie widać głównie Swift, a oferta to Flutter?}
\odp{Komercyjny Flutter i Fruit Tracker są prywatne. Publicznie pokazuję iOS i inne rzeczy. Na rozmowie mogę przejść po Flutterze, Firebase, Provider/Cubit i procesie sklepu z DEVEJI.}

\pyt{Mało testów w repo?}
\odp{Tak, w projektach uczelnianych/osobistych to słaby punkt. Wiem o tym i w komercyjnym kodzie chcę to robić lepiej.}

\pyt{Nie za junior?}
\odp{Mam komercyjny Flutter + sklepy + Firebase, własną apkę z ML/SQLite i mocne iOS. Szybko się uczę. Szukam miejsca, gdzie urosnę na realnych projektach.}

\section{Pytania behawioralne -- szkielet}

\pyt{Trudny bug?}
\odp{Sytuacja $\to$ co było do zrobienia $\to$ co zrobiłeś (reprodukcja, logi, fix) $\to$ efekt. Wypełnij prawdziwym przykładem z DEVEJI.}

\pyt{Deadline vs jakość?}
\odp{Tniesz zakres, nie fundamenty (crash, bezpieczeństwo). Trade-off mówisz wcześnie.}

\pyt{Z czego jesteś dumny?}
\odp{Wybierz jedno: publikacja w sklepie / Fruit Tracker / dopracowane powiadomienia w Water Reminder / architektura MacSpaceCleaner.}

\section{Gdyby jednak dali małe zadanie}

\pyt{Model z JSON?}
\odp{Klasa z polami + \texttt{fromJson}.}

\pyt{Lista z API?}
\odp{Stan loading $\to$ pobranie $\to$ success/error. UI: ListView.builder.}

\pyt{FutureBuilder czy Cubit?}
\odp{FutureBuilder na prostą, jednorazową rzecz. Cubit gdy stan wraca na wielu ekranach, chcesz retry i testy.}

\pyt{Szukajka z opóźnieniem?}
\odp{Poczekaj $\sim$300 ms po wpisaniu zanim wyślesz request. Anuluj poprzednie zapytanie.}

\pyt{Context po await?}
\odp{Sprawdź \texttt{context.mounted}. Inaczej możesz dostać błąd po wyjściu z ekranu.}

\section{Pytania, które Ty zadajesz Im}

\begin{enumerate}
  \item Jaki stan używacie na co dzień: BloC, Cubit czy Provider?
  \item Jak wygląda code review i wrzucanie do sklepów?
  \item Nad czym teraz pracuje zespół Flutter?
  \item Jak wygląda okres próbny / mentoring?
  \item Ile jest natywnego Swift/Kotlin obok Fluttera?
  \item Ile dni w biurze na start?
  \item Jakie są kolejne etapy rekrutacji?
\end{enumerate}

\section{Checklista przed rozmową}

\begin{itemize}
  \item 2--3 zdania o DEVEJI, Fruit Tracker, Water Reminder.
  \item Cubit vs BloC vs Provider -- umiesz powiedzieć własnymi słowami.
  \item Firebase Auth + Firestore -- krótko.
  \item Nie ściemniaj względem kodu (SwiftData, brak MapKit jeśli go nie ma).
  \item 3 pytania do nich.
  \item Adres: Nad Przyrwą 13, 10:00, ok.\ 45 minut.
\end{itemize}

\vspace{1em}
\noindent\textit{Mów prosto, konkretami z projektów. Jak czegoś nie wiesz -- powiedz jak byś to sprawdził. Powodzenia.}

\end{document}
